\newpage
\section{Cronograma de atividades do TCC}
\label{sc:cronograma}

As atividades planejadas para a execução deste Trabalho de Conclusão de Curso (TCC) estão distribuídas ao longo de cinco fases principais (F.1 a F.5), detalhadas na Tabela \ref{tab:fases}. A distribuição temporal destas atividades entre os meses de maio e setembro de 2026 está ilustrada nos gráficos de Gantt na sequência.

\vspace{0.5cm}

\begin{longtable}{p{\textwidth}}
\toprule%
\myrowcolour%
\bfseries F.1: Proposição e Delimitação \\
\midrule
Definição do tema, objetivos, pergunta de pesquisa e escopo do trabalho. \\
\midrule
\myrowcolour%
\bfseries Atividades planejadas para a fase 1 \\
\midrule
Elaboração do anteprojeto; reunião com orientador; definição da metodologia. \\
\midrule
\myrowcolour%
\bfseries Resultados esperados para a fase 1\\
\midrule
Anteprojeto aprovado pelo orientador e plano de trabalho definido.\\
\toprule%
\myrowcolour%
\bfseries F.2: Revisão Bibliográfica \\
\midrule
Levantamento e síntese da literatura relevante para fundamentar o TCC. \\
\midrule
\myrowcolour%
\bfseries Atividades planejadas para a fase 2 \\
\midrule
Pesquisa em bases; leitura crítica; organização das referências (BibTeX). \\
\midrule
\myrowcolour%
\bfseries Resultados esperados para a fase 2 \\
\midrule
Estado da arte consolidado e seção de revisão escrita.\\
\toprule%
\myrowcolour%
\bfseries F.3: Metodologia e Implementação \\
\midrule
Desenvolvimento do método, ferramentas e implementação do experimento/protótipo. \\
\midrule
\myrowcolour%
\bfseries Atividades planejadas para a fase 3 \\
\midrule
Desenho experimental; implementação do protótipo; testes iniciais. \\
\midrule
\myrowcolour%
\bfseries Resultados esperados para a fase 3 \\
\midrule
Protótipo funcional e conjunto de dados/resultados preliminares.\\
\toprule%
\myrowcolour%
\bfseries F.4: Análise, Redação e Entrega \\
\midrule
Análise dos resultados, redação dos capítulos e preparação da versão final. \\
\midrule
\myrowcolour%
\bfseries Atividades planejadas para a fase 4 \\
\midrule
Análise estatística/qualitativa; escrita dos capítulos; revisão com orientador. \\
\midrule
\myrowcolour%
\bfseries Resultados esperados para a fase 4 \\
\midrule
Versão final do relatório entregue e pronta para defesa.\\
\toprule%
\myrowcolour%
\bfseries F.5: Preparação da Apresentação/Defesa \\
\midrule
Preparação de slides, ensaio e ajustes finais para a defesa. \\
\midrule
\myrowcolour%
\bfseries Atividades planejadas para a fase 5 \\
\midrule
Criação de apresentação; ensaios; correções finais. \\
\midrule
\myrowcolour%
\bfseries Resultados esperados para a fase 5 \\
\midrule
Apresentação pronta e cronograma de defesa agendado.\\
\bottomrule
\caption{Detalhamento das fases do TCC}
\label{tab:fases}
\end{longtable}

\subsection*{Distribuição Temporal (2026)}

\noindent
% Configurações globais para os gráficos ficarem padronizados
\begin{ganttchart}[
  y unit title=.6cm,
  y unit chart=.7cm,
  x unit = \dimexpr\textwidth/6\relax, % Divide a largura da página perfeitamente em 6 colunas
  hgrid,
  vgrid={*1{draw=gray!30, dashed}}, % Linhas verticais tracejadas e suaves
  time slot format=isodate-yearmonth,
  time slot unit=month,
  inline, % Coloca o texto dentro da barra
  bar height=0.6,
  bar label font=\small\bfseries\color{white}, % Letra branca e negrito para contraste
  bar/.append style={fill=blue!60!cyan, rounded corners=2pt, draw=blue!80!black, thick} % Design moderno
]{2026-01}{2026-06}

  % Títulos do primeiro semestre
  \gantttitle{1º Semestre}{6} \\
  \gantttitle{Jan}{1}
  \gantttitle{Fev}{1}
  \gantttitle{Mar}{1}
  \gantttitle{Abr}{1}
  \gantttitle{Mai}{1}
  \gantttitle{Jun}{1} \\
  
  % Barras do primeiro semestre (Apenas a sigla para caber)
  \ganttbar{F.1}{2026-05}{2026-05} \\
  \ganttbar{F.2}{2026-05}{2026-06} \\
  \ganttbar{F.3}{2026-06}{2026-06} 
\end{ganttchart}

\vspace{1cm} % Espaço entre os dois semestres

\noindent
\begin{ganttchart}[
  y unit title=.6cm,
  y unit chart=.7cm,
  x unit = \dimexpr\textwidth/6\relax, 
  hgrid,
  vgrid={*1{draw=gray!30, dashed}},
  time slot format=isodate-yearmonth,
  time slot unit=month,
  inline,
  bar height=0.6,
  bar label font=\small\bfseries\color{white},
  bar/.append style={fill=blue!60!cyan, rounded corners=2pt, draw=blue!80!black, thick}
]{2026-07}{2026-12}

  % Títulos do segundo semestre
  \gantttitle{2º Semestre}{6} \\
  \gantttitle{Jul}{1}
  \gantttitle{Ago}{1}
  \gantttitle{Set}{1}
  \gantttitle{Out}{1}
  \gantttitle{Nov}{1}
  \gantttitle{Dez}{1} \\
  
  % Barras do segundo semestre
  \ganttbar{F.3}{2026-07}{2026-08} \\
  \ganttbar{F.4}{2026-08}{2026-09} \\
  \ganttbar{F.5}{2026-09}{2026-09}
\end{ganttchart}

\newpage